\documentclass[conference]{IEEEtran}
\usepackage{iftex}
\ifXeTeX
\usepackage{fontspec}
\setmainfont{Times New Roman}
\setsansfont{Arial}
\setmonofont{Courier New}
\else
\usepackage[T1]{fontenc}
\usepackage{lmodern}
\fi
\usepackage{amsmath,amssymb,array,url,graphicx,etoolbox}
\IfFileExists{booktabs.sty}{\usepackage{booktabs}}{\newcommand{\toprule}{\hline}\newcommand{\midrule}{\hline}\newcommand{\bottomrule}{\hline}}
\IfFileExists{kvsetkeys.sty}{\usepackage[hidelinks]{hyperref}}{}
\IfFileExists{microtype.sty}{\usepackage{microtype}}{}
\patchcmd{\abstract}{---}{: }{}{}
\patchcmd{\IEEEkeywords}{---}{: }{}{}
\title{LLM-Mediated Natural Language Controls for Short-Video Recommender Systems: An Offline Study of Controllability}
\author{\IEEEauthorblockN{Vibhor Malik}\IEEEauthorblockA{M.Sc. in Applied Computing\\British Columbia Institute of Technology\\Supervisor: Borna Noureddin\\COMP 9500 research project\\Review draft: 14 October 2026, Vancouver time}}
\begin{document}
\bstctlcite{controlNoAuthorDash}
\maketitle
\thispagestyle{plain}\pagestyle{plain}
\begin{abstract}
This offline study separates language translation from deterministic feed control. Frozen Experiment 1 increases target-category proportion from 11.6\% to 79.2\%, with muted exposure reduced to zero by hard exclusion; equivalent rules and Llama commands yield identical canonical rankings. The completed exposed-practice diagnostic retains its primary validated Llama score of 73/100 and its control-benefit/ranking-cost tradeoff. A prospectively frozen revision evaluates 100 fresh AI-authored requests against two separately contextualized AI references. Initial label agreement is 100/100 (Cohen kappa 1.0000); 0 disagreements receive prespecified AI adjudication. Validated parser agreement is 76/100 for Llama and 11/100 for rules, with 23 Llama failures retained. Induced control success is 20.14\% for Llama versus 3.82\% for item feedback; NDCG@10 is 0.02580 versus 0.03184. Full predictions and scores receive an independent computational audit. The references remain AI judgments, with shared model-family bias; these findings do not establish human validity, user satisfaction or general performance across LLMs. All 288 frozen Experiment 1 files remain intact.
\end{abstract}
\begin{IEEEkeywords}
controllable recommendation, natural-language interfaces, sequential recommendation, offline evaluation, reproducibility
\end{IEEEkeywords}

\section{Introduction}
Item feedback communicates a reaction to one recommendation. A request such as ``show more of this category'' specifies a broader intervention, while a profile can carry a preference across requests. Explanations serve a different purpose: making the information behind recommendations visible. Comparing these mechanisms requires separating the language model's interpretation from the ranker's mechanical enforcement.

Mozilla's 2022 audit of YouTube controls reported that Dislike prevented 12\% and Not interested 11\% of unwanted recommendations relative to its no-feedback control~\cite{ricks2022button}. This historical platform evidence motivates investigation; it is neither a current platform measurement nor a baseline for the present implementation. The research question here is narrower: how reliably does a local language model compile category requests, and how do category exposure and next-item accuracy change when the commands are applied to a frozen recommender?

The contribution is an inspectable applied artifact and evaluation pipeline. It does not introduce language-based recommendation or establish the novelty of combining controls. Parser validation, policy application, ranking metrics and explanations remain independently observable. The study uses historical public interactions and synthetic instructions, with no new participant recruitment. Consequently it cannot establish trust, naturalness, perceived agency or real-world satisfaction.

The supplied proposal is an unsigned draft, with unselected research-path boxes and blank approval fields. It specifies a 6--8-page IEEE or ACM paper; the supplied course-email transcript gives a broader 6--10-page range. Neither establishes whether references or appendices count. This main report includes its references within its page count and has no appendices. The preserved extended report and complete execution evidence are separate supplements; their existence does not establish supervisor approval of the implemented method.

\section{Related Work}
Language-based profiles are an established direction. Radlinski et al. outline transparent, scrutable profiles~\cite{radlinski2022profiles}; Wo\'{z}niak et al. investigate user-editable profiles~\cite{wozniak2025profiles}; Ramos et al. study edits and recommendation changes~\cite{ramos2024transparent}. These motivate explicit preferences, without showing that the present one-preference offline profile is equivalent to a rich human-edited profile. Sanner et al., Chen and Pu, and Zhou et al. provide related profile and critiquing approaches~\cite{sanner2023competitive,chen2012critiquing,zhou2024language}. Niu et al. examine conversational short-video control through interviews and a user study~\cite{niu2025chat}; Fabbri et al. describe an integrated transparent, controllable feed~\cite{fabbri2026feeding}. It would therefore be inaccurate to describe prior work universally as isolated single-channel systems.

Lu et al. align language models with recommendation instructions using a conversational recommender with a backbone such as SASRec~\cite{lu2024aligning}. Dai et al. study learned control tokens for reranking~\cite{dai2025token}. The current implementation instead freezes its recommender and uses the LLM only to translate commands; it cannot inherit those systems' reported benefits. SASRec supplies the architectural basis for the independently implemented compact network~\cite{kang2018sasrec}. KuaiRec provides dense historical short-video feedback and categories, not responses to the interventions constructed here~\cite{gao2022kuairec}.

Evaluation design is central. Ji et al. identify leakage risks when recommendation splits admit future interactions~\cite{ji2023leakage}. A category derived from a held-out item is still outcome-derived information even if the item's identifier is hidden. The revised pipeline prevents that path. Zhang et al.'s vision paper motivates preference-alignment measures beyond historical ranking accuracy~\cite{zhang2026context}; it does not validate this project's numerical satisfaction bands. Compliance, predictive relevance and human preference remain distinct constructs.

\section{Method}
\subsection{Data and information boundaries}
The processed cohort contains 865 eligible users, 2,593 candidate items and 31 opaque official category identifiers. The source KuaiRec matrix contains 1,411 users, 3,327 items and 4,676,570 interactions~\cite{gao2022kuairec}. A watch ratio above 2.0 is a derived positive-engagement proxy, not an observed like. Official documentation defines watch ratio as play duration divided by video duration~\cite{kuairecdocs}.

Preprocessing excludes 181,992 records without timestamps rather than inventing dates. Shared chronological boundaries divide source events into nominal 70\% training, 15\% request history and 15\% evaluation windows, before positive-event filtering and eligibility. Users need at least three training, two request-history and one evaluation item. Candidates must have a positive event before the training cutoff; 43,708 post-cutoff positive cold-item events are removed before user filtering. The retained histories contain 127,173 training positives, 11,712 request-history positives and 4,662 test positives. These counts describe different stages and must not be interpreted as one common-denominator retention fraction.

Training uses only training sequences. Scoring receives training plus request history, truncated to its maximum context. Requests may access that history and static categories, never final behavioral labels. The first eligible final item supplies the relevance label only after assignment and ranking. The frozen full-catalogue policy permits seen items, although the retained cohort has no repeated training items or test/history overlap; repeated targets therefore do not justify that policy. An unseen-only candidate design would be a separate experiment.

Categories remain identifiers such as \texttt{category\_12}; the study does not infer topic names or test understanding of natural video subjects. Synthetic requests are deterministic experimental assignments, not observations of what users wanted. Historical support does not make them independent human preferences or establish a causal response to an intervention.

\subsection{Recommender and matched configurations}
The compact SASRec-style network uses item and position embeddings, two self-attention blocks, one attention head, hidden dimension 32, maximum history 50 and dropout 0.2. It uses Adam with learning rate 0.001, batch size 64 and gradient clipping at 5.0, for 20 fixed epochs in seeds 42, 43 and 44. Negative sampling excludes training-positive items without consulting later positives. Checkpoints remain fixed across control conditions. The initial 256-user, three-epoch run measured execution cost before control-outcome inspection; the main run's fixed larger scope is retained in its provenance. No current run retrains or tunes the backbone.

Five configurations use the same catalogue and frozen score vector: A is simulated positive or negative feedback on one historical exemplar; B adds display-only explanations to A; C compiles a text command; D compiles one explicit profile preference; E combines text, profile and explanations without double-counting a repeated command. B must equal A in rank order. Equivalent C, D and E commands deliberately share policy. These configurations do not identify every mechanism's independent contribution: a category command communicates more than a single-item reaction, and explanations have no offline ranking intervention.

\begin{table}[t]\centering\small
\caption{Matched mechanisms and the comparison each supports}
\begin{tabular}{p{0.4cm}p{2.6cm}p{3.9cm}}\toprule
ID&Mechanism&Interpretation\\\midrule
A&One historical item reaction&Explicit simulated reference policy\\
B&A with explanation display&Ranking invariance, not a preference effect\\
C&Parsed free-text command&Translation followed by enforcement\\
D&One profile preference&Equivalent-intent consistency within session\\
E&Text, profile and explanation&Merge consistency without double counting\\\bottomrule
\end{tabular}\end{table}

For each user, category support comes from training and request history. A hash selects supported boost and mute targets; the latest matching historical item supplies A's exemplar. The free-text and profile channels express the same category operation. Their canonical wording was used during development, so that task is a consistency check rather than held-out language generalization.

\subsection{Parsing, enforcement and explanations}
The actual language backend is local Llama 3.1 8B through Ollama, with constrained structured output and application validation. A transparent rule parser is a separate comparator, never a substitute for failed inference. For E1, a command contains boost, mute, reset or clarify, a category and strength. Unknown categories, unsupported operations and malformed outputs remain errors. In end-to-end ranking, an error applies no new control and stays in the denominator.

Scores are normalized by each user's catalogue minimum and maximum, with tied scores mapped to zero. The default category boost adds 0.25; a hard mute removes matching items without padding forbidden items back into an underfilled list. Item identifiers break ties deterministically. Profile state is session scoped, and explicit commands take precedence over an equivalent stored preference. The profile does not implement the proposal's rich, approximately 200-word human-edited representation.

Explanations select one to three verbatim statements from a verified set of item, category and historical-association facts. This is grounded statement selection, not unrestricted explanation generation, causal faithfulness or evidence of human-perceived quality. Only ten baseline items have an actual model-generated cache; unmatched cards receive a visibly labeled verified template. Explanations do not change B's ranking.

The cache key binds the exact facts and model provenance. Changed history or item facts cannot legitimately reuse a mismatching explanation. Conversely, histories with the same verified facts may share a key. Ten cached cards therefore do not cover the full user/item space, and a visible template fallback must not be counted as another model inference. Source-code inspection or displaying a card is not evidence that a new explanation-generation experiment occurred.

Stateful interface checks are also separate from offline evidence. The API bounds inputs, isolates sessions/users and exposes backend/data status. Functional workflows cover setting preferences, changing users and resetting state. These checks can establish that the application applies a command consistently; they cannot determine whether a person understands a category identifier or prefers the resulting feed. A persistent profile across application restarts and a licensed video-playback experience are outside the implemented prototype.

\section{Evaluation}
\subsection{Metrics and statistical units}
At $K=10$, target-category proportion TCP@10 is the fraction of slots carrying a boosted target; TCER@10 is the same arithmetic for an unwanted target, where lower is better. The harness separately reports fill rate, conditional exposure for nonempty lists and filled compliance. Zero exposure after an empty output must not be mistaken for successful control. A positive presence check does not guarantee improved exposure over baseline.

Relevance uses one historical next-item label. HR@1 is one when it ranks first; NDCG@10 is $1/\log_2(r+1)$ when it appears at rank $r\leq10$, otherwise zero. These are next-item recovery proxies, not utility over all candidates or human satisfaction. A steering request can legitimately disagree with the historically observed next item, which predates the intervention.

E1 averages paired outcomes across requests and three seeds within each user. Thus 865 users remain 865 analysis units, rather than 5,190 independent observations. Its six primary contrasts are C, D and E versus A for TCP and TCER, with Wilcoxon signed-rank tests, Pratt zeros, rounded ties and Benjamini--Hochberg adjustment~\cite{scipywilcoxon,scipyfdr}. User-bootstrap intervals use 5,000 resamples. Shared configurations, sparse nonzero differences, ties and dependence restrict interpretation. Neither a small $p$ value nor a narrow interval corrects synthetic-intent bias or a weak backbone.

\subsection{Language accuracy and non-inferiority}
The frozen E1 broader suite measures exact operation/category/strength agreement with its fixed synthetic rubric. It maps ``less'' and ``fewer'' to hard mute, treats conditional requests as unsupported and asks for clarification when a neutral ``show'' request lacks direction. Those conventions are not established human meanings. One development/test pair differs only by a full stop; shared templates further limit independence. The separate 120 canonical inputs were development wording, and the ten adversarial cases form a small finite stress set.

The original parser target is 85\%, with a 70\% degradation floor. The broader suite's descriptive 95\% Wilson interval assumes a binomial working model; it does not establish population-wide reliability. Operational and validation failures remain in the denominator. Request/seed reuse cannot enlarge the number of unique language observations. A later prompt change would create development evidence and require a fresh prospective bank, not repair the interpretation of an already observed test.

The parser report distinguishes three levels: a model response may complete, its JSON may satisfy the supplied schema, and the validated command may agree with the reference. These are different outcomes. Application grounding can reject a syntactically valid command, while a valid accepted command can still express the wrong operation. End-to-end exact agreement treats both cases as failures. Reporting only completed HTTP requests or only cases that pass validation would inflate reliability. Raw responses, accepted commands, rejected commands and clarification outcomes must therefore remain traceable per question.

The proposal's ``within 10\%'' non-inferiority phrase is ambiguous. The frozen implementation used absolute HR@1 margin 0.10; the later audit additionally reports 10\% of the paired A baseline as a post hoc sensitivity analysis. With baseline about 0.0025048, that relative margin is about 0.00025048. The simultaneous lower bound for C/D/E minus A is about $-0.0007707$: above $-0.10$, below $-0.00025048$. The absolute calculation passes numerically but could permit complete baseline loss; the relative reading is unsupported. Neither establishes useful quality preservation or retroactive prespecification.

\section{Frozen Experiment 1 Results}
The executed cohort and checkpoints remain unchanged. Table~\ref{tab:e1} summarizes the matched control run. Target exposure changes are large, but policy construction explains their direction: category-wide boosting directly favors TCP against item-only feedback, and hard exclusion guarantees zero muted exposure after correct parsing. C, D and E have identical rankings on equivalent canonical commands; the separately executed rule comparator gives the same recorded rankings and aggregate metrics. This attributes the canonical-task gain to shared enforcement, not an observed LLM advantage.

\begin{table}[t]\centering\small
\caption{Frozen E1 control results, user means}\label{tab:e1}
\begin{tabular}{lrr}\toprule
Metric & A/B & C/D/E\\\midrule
Boost TCP@10 (\%) &11.615&79.222\\
Mute TCER@10 (\%) &6.112&0.000\\
Fill rate (\%) &100.000&100.000\\
Boost NDCG@10 &0.03012&0.01637\\
Mute NDCG@10 &0.03370&0.03526\\\bottomrule
\end{tabular}\end{table}

For C versus A, the mean paired boost TCP difference is 67.607 percentage points, with 95\% user-bootstrap interval [65.788, 69.438]; the TCER difference is $-6.112$ points, interval [$-7.056,-5.241$]. D and E yield the same contrasts. The BH-adjusted Wilcoxon values are approximately $1.35\times10^{-142}$ for boosts and $1.05\times10^{-65}$ for mutes. The proposal's directional control criterion is numerically supported under these conventions; this is not evidence that the LLM improves recommendation accuracy or understands general human intent.

Quality moves differently. Boost NDCG@10 decreases 45.7\%. The post hoc paired boost C--A difference is $-0.01375$, with user-bootstrap interval [$-0.02086,-0.00655$], and 108 users have nonzero differences. This secondary analysis was added after inspection and remains outside the unchanged primary BH family. Muted requests' NDCG increases, so one pooled metric would conceal distinct effects. The historical next-item label may disagree with a newly assigned steering instruction; the result still rules out an unqualified ``no quality cost'' claim.

The backbone is weak: per-seed HR@1 is 0.000\%, 0.347\% and 0.462\%, against about 7.977\% for request-window popularity. The latter uses pooled recent pre-test histories; the trained network uses each user's recent history for scoring but its weights were fitted on the earlier training window. Both avoid final labels, but this estimator difference prevents a controlled model-class superiority claim. Between 73.9\% and 96.2\% of users share the most-common first prediction within a seed. Such concentration and near-zero accuracy make non-inferiority especially uninformative.

\begin{table}[t]\centering\small
\caption{Historical uncontrolled backbone diagnosis}
\setlength{\tabcolsep}{3pt}
\begin{tabular}{lrrr}\toprule
Model&HR@1 (\%)&HR@10 (\%)&NDCG@10\\\midrule
Uniform, expected&0.039&0.386&0.00175\\
Training popularity&0.116&0.578&0.00298\\
Recent popularity&7.977&16.532&0.11866\\
SASRec seed 42&0.000&4.509&0.01642\\
SASRec seed 43&0.347&10.058&0.03627\\
SASRec seed 44&0.462&11.098&0.04584\\\bottomrule
\end{tabular}\end{table}

Llama parses all 120/120 canonical inputs, 79/100 broader synthetic requests and 2/10 adversarial requests. The broader-set Wilson interval is [70.02\%, 85.83\%]: the 70\% floor is cleared, while the 85\% point target is missed. The rule comparator obtains 81/100, not proof of general superiority from this finite set. Seven Llama cases are grounding rejections, recorded as errors rather than discarded; correct structured syntax and intended meaning are separate checks. All ten generated explanation cards equal the template, establishing no observed explanation advantage in this limited sample.

\begin{table}[t]\centering\small
\caption{Frozen E1 broader synthetic language results}
\begin{tabular}{lrr}\toprule
Case group&Rule correct&Llama correct\\\midrule
Positive&26/30&30/30\\
Negative&24/30&30/30\\
Reset&3/10&10/10\\
Ambiguous/unsupported&18/20&7/20\\
Adversarial&10/10&2/10\\\midrule
Total&81/100&79/100\\
Validation rejections&0&7\\\bottomrule
\end{tabular}\end{table}

The supported positive, negative and reset cases are all correct for Llama under the fixed E1 rubric. Difficulty concentrates in ambiguity and adversarial instructions. Fourteen broader-set cases yield valid but rubric-incorrect commands, in addition to the seven rejections. A rule can outperform a model on a narrow, explicitly coded rubric without demonstrating richer language understanding; a model can similarly clear a finite score floor without establishing safe interpretation of arbitrary requests. The result must be stated with its actual vocabulary, category scope and denominator.

Historical verification reconstructs prepared sequences and recomputes all 25,950 actual-model ranking rows. A later audit reproduces all 25,950 top-ten lists but reports 2,009 full-ranking hash differences and 5,190 request-score hash differences. Their cause remains unresolved; missing original score vectors prevent quantifying numerical error. Dependency pins alone do not guarantee cross-platform equality. These records are historical Linux checks, not a newly completed E1 replay on Windows or Colab.

\section{Preserved Exposed Practice Diagnostic}
The completed practice run at commit \texttt{c6ebda0} used 100 already exposed worked questions and AI-generated answers,25 in each of four intent classes. It is preserved unchanged and was not rerun. Primary validated whole-command agreement remains 73/100 for Llama and 18/100 for rules. Llama's 24 errors comprise 23 conditional-grounding rejections and one floor rejection; raw pre-validation agreement 95/100 is supplementary, never a replacement score. Conditional validated agreement is 2/25, compared with 25/25 compound,22/25 floor and 24/25 graded cases.

The practice run retains 3,459 assigned requests,41,508 ranking rows and 2,595 full score vectors. User-averaged induced success is 23.81\% for Llama,8.40\% for rules and 3.92\% for item feedback. Llama NDCG@10 is 0.02461 versus item feedback 0.03128, so the observed control benefit accompanies a predictive quality cost. Recency reaches 0.11866 NDCG@10. A separate standard-library recount reproduces all row metrics and pooled means without discrepancies; that earlier audit did not independently reimplement intervals or signed-rank tests. Exposed questions cannot support an unseen-test claim or enlarge the 100 unique language observations through request/seed reuse.

\section{Prospectively Frozen AI Reference Experiment}
\subsection{Design and provenance}
The user's revised authorization replaces the earlier human-gated experiment with study \texttt{AIREF-E2-20261014-v1}. The user reports oral approval from Borna; no written or independently verified supervisor receipt is asserted. The earlier protocol is retained as historical evidence. Human annotation is not a prerequisite of this revision and is absent from its validity claims.

An independent AI reviewer checked the prospective protocol before the source/design freeze at 20:44:17 UTC on 14 October 2026. The protocol fixes the research question, sampling, prompts, rubric, failures, separation, metrics and disagreement handling before any fresh wording. Vibhor is an AI custodian; a separate clean author context receives only the sampling plan and authoring brief. A seeded 100-slot plan balances compound, graded, nonzero-floor and same-item conditional requests (25 each), covering all 31 target identifiers. The grade mix is nine slight boosts, eight strong reductions and eight exclusions. The author supplies three wording candidates per slot. First passing candidates are selected in fixed order with no parser predictions.

Screening excludes exact normalized duplicates and close lexical matches to the exposed 100 questions and 48 archived development texts. The prespecified flags are normalized SequenceMatcher similarity at least 0.82 or token-bigram Jaccard at least 0.55, after category placeholders and punctuation normalization. All candidates and exclusions remain saved; fresh within-bank near-paraphrase flags are reported, not selected on outcomes. These heuristics do not guarantee semantic or pretraining novelty.

Tanishq and Khushi are AI reference annotators in separate fresh contexts. They receive independently shuffled exact question sheets and the fixed semantic rubric, without author plans, implementation, predictions or each other's labels. Their two original label files and receipts are locked before comparison. Exact agreements become final; a fresh AI adjudicator resolves disagreements without predictions. No original annotation is overwritten. Final references must satisfy the sampled subtype and target roles; incompatible or unresolved decisions stop inference rather than being coerced. All roles inherit GPT-6.1-sol with xhigh reasoning, with actual task/input declarations. Exact serving weights are unavailable; separate contexts share a model family and may share errors. Their labels are AI references rather than authenticated human truth.

The first reference-custody attempt aborted before comparison/inference on a receipt-field mismatch. Its originals remain archived. A separately frozen second attempt reused the identical selected bank and unchanged semantic protocol; fresh, separately contextualized annotators wrote both initial label sets anew. The administrative amendment and independent pre-annotation review preserve this chronology and prevent the failed receipts from being silently repaired.

\subsection{Execution and metrics}
All final cases, both labels, adjudication, sources, prompts, validators, ranking inputs, model digest, seeds and metrics receive a full pre-inference hash lock. The Llama runner receives only opaque IDs, exact requests and the 31 public category IDs. Trusted worker code applies file/network audit restrictions and verifies denied reads of the protected label files; no reference contents enter Ollama's messages. Allowed networking is loopback 11434 only. This is application separation for trusted code on a shared Windows host, not an OS security sandbox against a compromised administrator.

The unchanged local \texttt{llama3.1:8b} / \texttt{Q4\_K\_M} model has SHA-256 digest \texttt{46e0c10c039e0191}\allowbreak\texttt{19339687c3c1757c}\allowbreak\texttt{c81b9da49709a3b3}\allowbreak\texttt{924863ba87ca666e}, using Ollama 0.34.0. Fixed options are temperature 0, seed 42, context 4096, output limit 768, six threads and 600-second timeout. Rules and Llama each attempt every exact request once. Durable pre-call journals and terminal records prevent silent reruns; any interruption aborts the version. Failures apply no new ranking control and remain in parser/ranking denominators. All 200 terminal predictions are committed before protected scoring joins references. No prompt, validator, band or threshold is tuned from practice or new outcomes.

The same 865-user,2,593-item data and three frozen SASRec checkpoints (42/43/44) supply matched item-feedback, rules, Llama and recency conditions. Eligibility depends on pre-test histories and final AI reference categories; SHA-256 assignment seed 2026 selects one eligible request per user/class. Next-item labels are accessed only for relevance after rankings. No retraining occurs. Induced success requires a full top-ten list, the frozen operational count band and positive signed movement $M>10^{-12}$. Conditional movement penalizes added nonintersection exposure; floors require nonzero presence. These numerical bands are conventions rather than human satisfaction judgments.

Users are analysis units, averaging three seeds within each request/class and eligible classes equally within user. Paired 5,000-resample percentile bootstrap intervals use seed 2026. Eleven one-sided Pratt Wilcoxon comparisons use 12-decimal rounding, asymptotic ties and joint BH correction: four classes plus pooled Llama--A and Llama--rules success, and pooled Llama--A movement. H1\_AI needs both positive significant pooled success contrasts; H2\_AI needs positive significant movement. HR@1 and NDCG@10 are reported together with control outcomes as descriptive H3\_AI, without a noninferiority claim. Complete condition/seed masks preserve unfavorable outcomes, infeasibility, missing opportunity and exclusions.

\subsection{Results}
Initial whole-command AI agreement is 100/100, Cohen kappa 1.000000; 0 disputes received the prescribed AI adjudication. Directive and field agreement, original labels, rationales and decisions remain in the supplement. This measures agreement among AI contexts, not factual correctness.

\begin{table}[t]\centering\small
\caption{Fresh AI-reference parser results; failures retained}\label{tab:aiparser}
\begin{tabular}{lrrr}\toprule
Class & Rules & Llama & Llama errors\\\midrule
Compound & 6/25 & 25/25 & 0\\
Graded & 5/25 & 24/25 & 0\\
Floor & 0/25 & 25/25 & 0\\
Conditional & 0/25 & 2/25 & 23\\
All & 11/100 & 76/100 & 23\\\bottomrule
\end{tabular}\end{table}

The run preserves 3,460 assigned user/class requests,41,520 ranking rows and 2,595 full score vectors. There are 0 eligibility exclusions. All 23 Llama failures are unchanged-validator rejections of condition grounding. Parser errors remain in 100-case denominators and every assigned ranking comparison.

\begin{table}[t]\centering\small
\caption{Fresh AI-reference ranking outcomes, paired user means}\label{tab:airank}
\begin{tabular}{lrrr}\toprule
Condition & Success(\%) & HR@1(\%) & NDCG@10\\\midrule
Item feedback & 3.815 & 0.241 & 0.03184\\
Rules & 3.854 & 0.231 & 0.03097\\
Llama & 20.145 & 0.337 & 0.02580\\
Recency & 5.202 & 7.977 & 0.11866\\
\bottomrule
\end{tabular}\end{table}

Under this revision only, H1\_AI is supported and H2\_AI is supported by the prespecified directional tests with joint BH correction across 11 contrasts. H3\_AI reports predictive quality descriptively. These decisions do not complete the archived human-reference design.

User-averaged induced control success is 20.145\% for Llama, 3.854\% for rules and 3.815\% for item feedback. NDCG@10 is 0.02580 for Llama versus 0.03184 for item feedback. Control and next-item ranking quality must be interpreted together.

A separate AI reviewer audits immutable evidence and independently recalculates the parser and ranking results; the computational audit records its exact statistical scope. No reference is supplied to inference during review. All 288 E1 files remain byte-identical. Reused questions across users and seeds do not increase the 100 distinct language observations.

\section{Validity and Reproducibility}
Synthetic opaque categories restrict construct validity. A single historical exemplar conveys less information than a category instruction. Equivalent policy outputs do not establish equivalent usability, while hard exclusion's zero exposure is true by construction. The study neither observes real steering intentions nor measures human comprehension or satisfaction. KuaiRec's platform, period, positive-event threshold and filtered active cohort restrict external validity.

Shared templates and exposure limit language independence. The weak untuned backbone, sparse paired quality differences and globally pooled recency comparator complicate quality interpretation. The original absolute and relative non-inferiority readings remain visible. No post hoc approval can make a margin prospectively specified, an exposed bank unseen, or a computational audit human annotation. Unresolved historical ranking/score hashes remain disclosed alongside successful top-ten reproduction.

The artifact separates preprocessing, training, controls, evaluation and a React/FastAPI demonstration. Real-data versus fixture status and rule versus local-model status are visible. Metadata cards are not playable video assets. State is session scoped; malformed commands remain errors. Prior unit/API/browser checks support software behavior, not the research hypotheses or present full-platform acceptance. The Windows practice and prospective AI-reference runs are verified; full E1 replay and Colab execution remain unverified.

The separate methods bundle retains the literal prompt/schema, runtime digest, options, source snapshots, input hashes, original reference bytes, 100 raw responses, 200 parse records, score vectors, ranking rows, exclusions, failures, terminal logs and audit evidence. It also preserves the interrupted wrong-model attempt, from which no result is used. New complete score vectors improve traceability for this run without resolving the older E1 discrepancy. Dataset/model licenses and pseudonymous identities must be respected; model weights are not redistributed.

The local execution environment is identified by a real terminal receipt from AI\_WORK, not inferred from a cloud workspace. The installed Ollama service responded on loopback and its model manifest digest matched the archived Llama identity after the authorized download. Python was 3.14.0, with NumPy 2.3.5, SciPy 1.16.3 and PyTorch 2.9.1 CPU; the exact environment record is preserved. This differs from historical package snapshots. Passing byte hashes proves artifact identity, not numerical equivalence across all runtime versions.

Reproduction should start from those saved artifacts: authenticate the E1 inventory, verify current source/input hashes and runtime identity, then replay the cached parsing and recompute rankings from the recorded scores. A fresh inference call is a new execution and should receive a new directory and timestamp; it must not silently replace these completed records. The public repository remains a partial review copy, while the full local bundle retains pseudonymous request-level evidence and scores. The compiled main report and separate supplementary report do not replace the experiment's frozen inputs.

The project used substantial AI assistance for implementation, analysis and drafting. Responsibility for understanding, source checking, policy compliance and genuine activity records remains with the student. This paper claims no new student hours, personal reading, meeting occurrence or independent human annotation from automated execution.

This is a prospectively frozen evaluation of fresh synthetic language against AI references. It does not establish independent human validation, naturalistic user benefit, contamination-free pretraining, general performance across LLMs, or guaranteed recommendation accuracy. The weak backbone, opaque categories, finite bank, historical relevance proxy and correlated annotator family constrain interpretation.

\section{Conclusion}
Frozen E1 demonstrates bounded category enforcement and predictive quality cost, with no canonical-task LLM ranking advantage. The exposed practice result remains a separate 73/100 validated diagnostic. In the fresh, prospectively frozen AI-reference revision, Llama agrees on 76/100 whole commands versus 11/100 for rules, with 23 retained validator failures. Its user-averaged induced control success is 20.145\% versus 3.815\% for item feedback, while NDCG@10 falls to 0.02580 versus 0.03184. H1\_AI and H2\_AI meet the prespecified directional tests for this fixed AI-reference design. Shared-model AI labels do not establish independent human truth, naturalness, satisfaction or performance across LLMs. Full score vectors, original labels and failures make the bounded result auditable; human experience remains a separately designed extension.

\bibliographystyle{IEEEtran}
\bibliography{references}
\end{document}
